\section{Leakage in Cluster-Scale KVCache}
\label{sec:attacks}

We first ask what an API tenant can learn from a \mc-style deployment, under
both unrestricted sharing and per-principal salting. \Cref{tab:attacks}
summarizes the six attacks. \A1 and \A5 adapt known single-engine
channels~\cite{song2024earlybird,cve2026105752} to the cluster; \A2, \A3,
\A4 and \A6 arise from \mc's global pool, tiers, replication and
admission control.

\begin{table*}[t]
\centering\small
\begin{tabular}{@{}l l l l c c c@{}}
\toprule
 & Channel & Mechanism & Attacker learns & Timing? & Survives salting? & Result \\
\midrule
\A1 & Global hit & Cross-node block fetch beats recompute & Membership & yes & no & \tbd{AUC} \\
\A2 & Tier ladder & HBM / local DRAM / remote DRAM / SSD latency classes & Recency & yes & no & \tbd{AUC} \\
\A3 & Replication & Replica count changes fetch contention & Popularity & yes & no & \tbd{AUC} \\
\A4 & Admission & Hit admitted, miss gets HTTP 429 near the SLO edge & Membership & \textbf{no} & no & \tbd{AUC} \\
\A5 & Accounting & Usage record reports cached tokens & Exact hit length & \textbf{no} & no$^\dagger$ & \tbd{AUC} \\
\A6 & Eviction & Victim traffic evicts attacker-primed blocks & Activity volume & yes & \textbf{yes} & \tbd{AUC} \\
\bottomrule
\end{tabular}
\caption{Attacks on a \mc-style deployment. ``Survives salting'' means the
channel persists when every principal has its own cache key namespace.
$^\dagger$Unless salting is bypassed, as in~\cite{cve2026105752}.
Results: \tbd{distinguishability (AUC) at moderate load, \cref{sec:eval-attacks}}.}
\label{tab:attacks}
\end{table*}

\subsection{Setup and method}
\label{sec:attack-method}
\todo{Finalize after week-1 testbed. Planned: \tbd{2} nodes $\times$ 8 \tbd{GPU
model}, \tbd{2$\times$200\,Gbps} RoCE, vLLM \tbd{version} with the \mc Store
connector, a Conductor-style router implementing Algorithm~1
of~\cite{qin2025mooncake}, models \tbd{Llama-3.1-8B, Qwen 14B, 70B-class}.}
Background traffic replays the conversation and tool\&agent traces released
with \mc~\cite{qin2025mooncake} at \tbd{three} load levels. Victims send prompts
from templated sensitive corpora (system prompts; prompts embedding names and
medical or financial attributes, following~\cite{chu2025safekv,pennas2026cachesolidarity}).
For each attack we train a classifier on attacker-only calibration traffic and
report AUC, attack success rate and queries per recovered token. We also measure
how load erodes each channel, since contention can make quiet-system
measurements optimistic~\cite{abubakar2026contention}.

\subsection{Timing channels}

\paragraph{\A1: Global hit.} The attacker sends a candidate prefix. If any
principal's request placed matching blocks anywhere in the pool, the Conductor
routes the request to fetch them and TTFT drops by the saved prefill minus the
RDMA transfer. Unlike single-engine attacks, the victim and attacker need not
share a node. \tbd{Result: AUC, success rate per token, effect of prefix length.}

\paragraph{\A2: Tier ladder.} A hit is not binary. A block may sit in the
serving node's HBM or DRAM, in another node's DRAM, or on SSD, each with a
distinct load time; which tier holds it depends on how recently and where it
was last used. Probing a known prefix and classifying the latency class reveals
\emph{when} it was last touched, not only \emph{whether}. \tbd{Result: per-tier
TTFT distributions (\cref{fig:ladder}), classification accuracy.}

\paragraph{\A3: Replication.} \mc replicates blocks that many requests
fetch~\cite{qin2025mooncake}. More replicas mean less contention on any one
holder, so repeated probes of a popular prefix see lower and less variable
fetch times. The attacker estimates how many principals use a prefix---for
example, how much traffic a competitor's agent receives. \tbd{Result: correlation
between estimated and true request rate.}

\begin{figure}[t]
\centering
\placeholderfig{1.55in}{TTFT CDFs for a 4k-token prefix: miss vs.\ hit from
local HBM, local DRAM, remote DRAM, SSD, at low and moderate load.}
\caption{\A2: the tier ladder. \tbd{Takeaway sentence with numbers.}}
\label{fig:ladder}
\end{figure}

\subsection{Non-timing channels}

\paragraph{\A4: Admission.} The Conductor rejects a request when predicted
TTFT exceeds the SLO, and the prediction uses the global prefix match. The
attacker raises its prompt length until, at the current load, a miss would be
rejected and a hit would not; HTTP 429 versus 200 then answers the membership
question with no clock at all, immune to network jitter. \tbd{Result: success
rate versus load, number of calibration queries.}

\paragraph{\A5: Accounting.} Usage records that report cached tokens give the
exact hit length per request. With global sharing this is a perfect membership
oracle; with salting it leaks whenever the salt is lost along any code
path~\cite{cve2026105752}. \tbd{Result: confirm on testbed; survey of what
public APIs report, with disclosure (\cref{app:ethics}).}

\subsection{A contention channel that survives salting}

\paragraph{\A6: Eviction.} Salting separates key spaces but not capacity. The
attacker primes long, unique prefixes of its own, waits, and re-probes them.
Global LRU evicts the attacker's blocks only when others insert enough new
blocks, so the fraction of the attacker's prefixes that survive measures
other principals' cache insertion volume over the interval, with time
resolution set by the probe rate. This is a Prime+Probe analogue at the
storage tier~\cite{yarom2014flush}. \tbd{Result: correlation between surviving
fraction and victim insertion rate; salted vs.\ unsalted.}

\subsection{Takeaways}
\todo{Write once results exist. Expected shape: (1) the global pool widens
membership attacks across nodes; (2) the tiers turn membership into recency and
popularity; (3) admission and accounting leak with no timing, so timing-only
defenses such as padding cannot close them; (4) salting removes \A1--\A5 at the
price of all cross-principal reuse, and still leaves \A6.}
